\section{The two-body kernels of $Q_1$, $S_1$ and $K_1$}

For $x$ in a homogeneous basis of $\mathfrak{psl}(4|4)$, $\lambda_x$ is its $\mathrm{ad}_\Delta$ eigenvalue, $-1$ on $\Sigma$, $+1$ on $\Sigma^*$ and $+2$ for $K_b$, and the degree of $x$ in the position is $0$, $1$ and $2$ (Section 10.9). The two-body kernel of a pair of letters along $x$ is that of Proposition \ref{prop:c10-arity1}: the limit as $\varepsilon\to0$ of the local parts of the cycles through the pair with exactly one defect insertion, a special edge, the vertex $S^{(1)}(x)$ on the cycle, or the arity-one counterterm vertex of the operator. This section gives the lines that $S^{(1)}(x)$ forms on a cycle along $\Sigma$, $\Sigma^*$ and $K_b$, the output bounds, the relations of $S_1$ to $Q_1$ and of $K_1$ to $P_1$, the renormalisation freedom in the scheme of Sections 11.3--11.5, and the comparison of $Q_1$ with the one-loop differential of the holomorphic twist.

\begin{lemma}[The lines along $\Sigma$, $\Sigma^*$ and $K_b$]
\label{lem:c11-lines}
\begin{itemize}
\item (i) For $\xi\in\Sigma$ the quadratic part of $S^{(1)}(\xi)$ is $\int\mathrm{dvol}\big(-(\Gamma(\xi,\lambda),A^*)+\tfrac12(\rho(dA)\xi,\lambda^*)\big)$ (Proposition \ref{prop:c7-susy}). With both legs on edges $P$, its term $-\int(\Gamma(\xi,\lambda)|_V,A^*|_V)$ has the leg $\lambda$ on the component $\lambda\otimes\lambda$ of $P$ and the leg $A^*|_V$ on the component $A^*|_V\otimes c$: an odd line joining a slot $\lambda$ to a slot $c$, linear in $\xi$ and free of the position. The terms $-\int(\Gamma(\xi,\lambda)|_{W_{\mathbb R}^\perp},A^*|_{W_{\mathbb R}^\perp})$ and $\tfrac12\int(\rho(dA)\xi,\lambda^*)$, with a leg $A^*|_{W_{\mathbb R}^\perp}$ or $\lambda^*$, enter a cycle at the special edge.
\item (ii) For $\eta\in\Sigma^*$, $S^{(1)}(\eta)$ is $S^{(1)}(\xi)$ with $\xi$ replaced by the value $\rho(x)\eta$ of the conformal Killing spinor, plus the translation-invariant term $-\int(\rho(A|_{W_{\mathbb R}^\perp})\eta,\lambda^*)$ \eqref{eq:c2-S1S}. Its line on a cycle is the line of (i) along $\rho(x)\eta$, linear in the position; the last term, with its leg $\lambda^*$, enters a cycle at the special edge.
\item (iii) For $K_b$, with the vector field $v(x)=(x,x)b-2(b,x)x$, $S^{(1)}(K_b)$ forms on a cycle the composite line from $c$ to $A|_V$ with the kernel $P\star\,v.\,P\star=v.(P\star)^2+[P\star,v.]P\star$, quadratic in the position, at a finite endpoint $T$ (Definition \ref{def:c8-K} and the remark after Theorem \ref{thm:c8-K}).
\end{itemize}
In each case every letter, $A^*|_{W_{\mathbb R}^\perp}$, $\lambda^*$ and $c^*$ included, is active (Lemma \ref{lem:c10-active1}).
\end{lemma}

\begin{proof}
The components of $P$ are those of Proposition \ref{prop:c3-P}; $\Gamma(\xi,\lambda)\in V_{10}$ pairs with $A^*$ through its parts in $V$ and in $W_{\mathbb R}^\perp$. (ii) is \eqref{eq:c2-S1S}, the first two terms being those of \eqref{eq:c2-S1Q} at the spinor $\rho(x)\eta$. (iii) Definition \ref{def:c7-vdot} for the vector field of $K_b$ and Lemma \ref{lem:c11-kernels} (iii), (iv) with this $v$.
\end{proof}

\begin{proposition}[Output bounds along $\Sigma$, $\Sigma^*$ and $K_b$]
\label{prop:c11-oddbound}
Under (iii) of Definition \ref{def:c10-hyp} along $x$, the two-body kernel of $(X_a,X_b)$ is a local functional of the pair degree with:
\begin{itemize}
\item (i) along $\xi\in\Sigma$: Weyl weight $w_{X_a}+w_{X_b}+\tfrac12$, no position, parity $|X_a|+|X_b|+1$, and at most $w_{X_a}+w_{X_b}+\tfrac12$ plus the pair degree output fields, as for the second-order operator on pairs of letters of [BK, (3.7)];
\item (ii) along $\eta\in\Sigma^*$: Weyl weight $w_{X_a}+w_{X_b}-\tfrac12$ with the position at weight $-1$, at most linear in the position, parity $|X_a|+|X_b|+1$, and at most $w_{X_a}+w_{X_b}+\tfrac12$ plus the pair degree output fields;
\item (iii) along $K_b$: Weyl weight $w_{X_a}+w_{X_b}-1$ with the position at weight $-1$, at most quadratic in the position, parity $|X_a|+|X_b|$, and at most $w_{X_a}+w_{X_b}+1$ plus the pair degree output fields.
\end{itemize}
\end{proposition}

\begin{proof}
Proposition \ref{prop:c10-bound1} with $\lambda_x=-1,+1,+2$ and the degrees $0,1,2$ in the position.
\end{proof}

\begin{proposition}[The special supercharges and the supercharges]
\label{prop:c11-S}
Let $\eta\in\Sigma^*$ and $b\in V$, so that $\rho(b)\eta\in\Sigma$.
\begin{itemize}
\item (i) Graph by graph, the position part along $b$ of the two-body kernel of $(X_a,X_b)$ along $\eta$ is the two-body kernel of $(X_a,X_b)$ along $\rho(b)\eta$; its translation-invariant part is computed with $S^{(1)}(\eta)$ at the spinor $\rho(x-x_O)\eta$, $x_O$ the point of the density, together with the term $-\int(\rho(A|_{W_{\mathbb R}^\perp})\eta,\lambda^*)$.
\item (ii) Under the hypotheses of Proposition \ref{prop:c7-select} and Corollary \ref{cor:c8-homotopy}, on cohomology, the position part along $b$ of $S_1$ along $\eta$ is $Q_1$ along $\rho(b)\eta$ minus the commutator of $\{S^{(1)}(\eta),-\}$ with $P_1$.
\item (iii) On a family on which $P_1\cong0$, in particular on the families of Definition \ref{def:c11-families} in the translation-covariant scheme at $\theta=\theta_\lambda=\theta_V=-\tfrac13$, the position part of $S_1$ along $\eta$ is $Q_1$ along $\rho(b)\eta$ on cohomology, and $S_1$ along $\eta$ is fixed by $Q_1$ along $\Sigma$ up to the term of Proposition \ref{prop:c7-select} linear in ${\not d}(\rho(x)\eta)=4\eta$. On the scalar bilinears this term is absent (Corollary \ref{cor:c7-S}), and $S_1$ along $\eta$ is, on cohomology, the combination of the functionals of Corollary \ref{cor:c7-S} with the coefficients of $Q_1$ along $\Sigma$ in Corollary \ref{cor:c7-Q}.
\end{itemize}
\end{proposition}

\begin{proof}
(i) By Lemma \ref{lem:c11-lines} (ii), $S^{(1)}(\eta)=\int x^\mu s_\mu+\int s_0$, where $\int s_\mu$ is $S^{(1)}$ along the constant spinor $\rho(e_\mu)\eta\in\Sigma$ and $s_0=-(\rho(A|_{W_{\mathbb R}^\perp})\eta,\lambda^*)$. Insert $x^\mu=x_O^\mu+(x-x_O)^\mu$ and argue as in the proof of Proposition \ref{prop:c11-position} (i): on plane waves the factor $x^\mu$ is $-i\partial/\partial r^\mu$ of the momentum injected at the vertex, and its action on the momentum delta function is $-i$ times the graph with $\int s_\mu$; the same holds for the special edges, the vertices $S^{(1)}(\eta)$ on the cycle and the counterterm vertices, which are the defect insertions of Proposition \ref{prop:c10-arity1}. (ii) Corollary \ref{cor:c8-homotopy} (ii) at $(\eta,P_b)$: by Proposition \ref{prop:c2-special}, $[P_b,\eta]=\rho(b)\eta$, so $[\eta,P_b]=-\rho(b)\eta$, and $S^{(2)}(\eta,P_b)=0$ (Definition \ref{def:c2-S2}). At order $\hbar$ the identity states that the commutator of $S_1$ along $\eta$ with $\{S^{(1)}(P_b),-\}$, plus the commutator of $\{S^{(1)}(\eta),-\}$ with $P_1$, minus $Q_1$ along $\rho(b)\eta$, is a graded commutator with $\Lambda^{(0)}$, which induces $0$ on cohomology; at order $\hbar^0$ the same identity is the statement of (i) for the classical transformation, which fixes the signs. By Proposition \ref{prop:c10-bound1}, $S_1$ along $\eta$ is at most linear in the position, and Lemma \ref{lem:c11-canonical} identifies the first commutator with its position part. (iii) By Proposition \ref{prop:c7-select} the representative of $S_1$ along $\eta$ at a point $x_0$ is a term linear in the value $\rho(x_0)\eta$ plus a term linear in $4\eta$; the first has the position part given by its value at $\rho(b)\eta$, and as $\eta$ and $b$ vary, $\rho(b)\eta$ spans $\Sigma$, so by (ii) it is the representative of $Q_1$ evaluated at $\rho(x_0)\eta$ on cohomology.
\end{proof}

\begin{proposition}[Renormalisation freedom of the odd components and of $K_1$]
\label{prop:c11-freedom}
\begin{itemize}
\item (i) A field redefinition, a closed direction and the outer direction $G_B$ (Proposition \ref{prop:c8-laws} (a), (c)) change each component by minus the bracket with a local functional, a derivation of the letters: they change the one-body kernels and leave the two-body kernels unchanged. $G_B$ changes $Q_1$ along $S_+\otimes W$ by $-a$ and along $S_-\otimes W^*$ by $+a$ times the classical transformation, $S_1$ along $S_+\otimes W^*$ by $+a$ and along $S_-\otimes W$ by $-a$ times it, and leaves $K_1$ unchanged.
\item (ii) The finite counterterms $\hbar F_1^{(0)}$ and $\hbar F_1^{(1)}(x)$ enter the one-body kernels (Remark \ref{rem:c10-scheme1} (a), (b)); the sector coefficients of Definition \ref{def:c11-free} enter the components along $\Delta$ and $P_b$.
\item (iii) A finite part $\psi'$ of the operators changes the component along $x$ by $[\{S^{(1)}(x),-\},\psi']$ (Proposition \ref{prop:c8-laws} (b)). For $\psi'=\tfrac a2\{S^{(1)}(\Delta),-\}$ this is $\tfrac a2\lambda_x\{S^{(1)}(x),-\}$ (Proposition \ref{prop:c9-scheme}): $+a$, $-\tfrac a2$, $+\tfrac a2$ times the classical transformation along $K_b$, $\Sigma$, $\Sigma^*$, a change of the one-body kernels. A finite part $\psi'$ that vanishes on $O$ changes the component of $O$ along $x$ by $-\psi'(\{S^{(1)}(x),O\})$; in this way the finite parts of the operators of the Weyl weight of $\{S^{(1)}(x),O\}$ enter the two-body kernels.
\item (iv) The cutoff shape enters the finite parts of the special-edge cycles (Proposition \ref{prop:c10-shape}, Lemma \ref{lem:c10-special}), the composite pair kernels and the finite vertex \eqref{eq:c10-finitesp}.
\item (v) In the translation-covariant scheme at $\theta=\theta_\lambda=\theta_V=-\tfrac13$, $P_1\cong0$ on the families of Definition \ref{def:c11-families} (Corollary \ref{cor:c11-D1fam}), and the condition $P_1=0$ of Definition \ref{def:c8-radial} holds there with $\psi'=0$: the components along $K_b$, $\Sigma$ and $\Sigma^*$ are those of the scheme. For equal coefficients $\theta_P\ne\theta$ along $P_b$, the finite part with $a=(\theta-\theta_P)$ units gives $P_1\cong0$ and changes the components along $K_b$, $\Sigma$, $\Sigma^*$ by $a$, $-\tfrac a2$, $\tfrac a2$ times the classical transformation (Corollary \ref{cor:c11-theta} (iv)).
\item (vi) The class $[Q^\xi_1]$ on the cohomology of $Q^\xi_0$, its kernel, image and rank are invariant under (i) and (iii) (Theorem \ref{thm:c8-twist}).
\end{itemize}
\end{proposition}

\begin{proof}
(i) Proposition \ref{prop:c8-laws}; the bracket with a local functional is a derivation in its second argument, so it acts letter by letter. (ii) Remark \ref{rem:c10-scheme1} and Definition \ref{def:c8-components}. (iii) Proposition \ref{prop:c8-laws} (b) and Proposition \ref{prop:c9-scheme}, with $[\{S^{(1)}(x),-\},\psi'](O)=\{S^{(1)}(x),\psi'(O)\}-\psi'(\{S^{(1)}(x),O\})$. (iv) Remark \ref{rem:c10-scheme1} (d). (v) Corollaries \ref{cor:c11-D1fam} and \ref{cor:c11-theta} (iv). (vi) Theorem \ref{thm:c8-twist}.
\end{proof}

\begin{corollary}[$K_1$ in the translation-covariant scheme]
\label{cor:c11-K}
Under the hypotheses of Theorem \ref{thm:c8-adjoint}, on a family of Definition \ref{def:c11-families} at the origin of one length, in the translation-covariant scheme at $\theta=\theta_\lambda=\theta_V=-\tfrac13$, $K_1$ along $K_{\bar b}$ equals $\hbar\big[\{S^{(1)}(K_{\bar b}),-\},\mathrm{Gram}_{\rm tree}^{-1}\mathrm{Gram}_{\rm 1}\big]$, the statement $K=P^\dagger$ of Corollary \ref{cor:c8-gram}. For equal coefficients $\theta_P$ along $P_b$, $K_1$ along $K_{\bar b}$ is $(\theta-\theta_P)$ units times $\{S^{(1)}(K_{\bar b}),-\}$ plus the same commutator, with the radial Gram matrix of Definition \ref{def:c8-inner} in that scheme.
\end{corollary}

\begin{proof}
Corollary \ref{cor:c11-D1fam} gives $P_1\cong0$, resp. $P_1\cong(\theta-\theta_P)$ units times $\{S^{(1)}(P_b),-\}$ with a real coefficient; Corollary \ref{cor:c8-gram} gives the two relations.
\end{proof}

\begin{proposition}[The chiral-primary bilinears along a rank-one supercharge]
\label{prop:c11-plane}
Let $\xi\in S_+\otimes W$ be of rank one, $\xi_\alpha{}^i=u_\alpha z^i$, and let $w_1,w_2\in W_{\mathbb R}^\perp\otimes\mathbb C$ satisfy $\rho(w_1)\xi=\rho(w_2)\xi=0$. Then $(w_1,w_2)=0$, and under the hypotheses of Proposition \ref{prop:c7-select} the component along $\xi$ of $O=\int\omega\,(\iota_{w_1}A,\iota_{w_2}A)$ vanishes at every order at which they hold: its classical part $\{S^{(1)}(\xi),O\}$ and $Q_1(O)$ are $0$, in every scheme in which the counterterms are equivariant for translations, $\mathfrak{so}(V)$, $\mathfrak{so}(W_{\mathbb R}^\perp)$ and $\mathfrak g$.
\end{proposition}

\begin{proof}
$\rho(w_1)\rho(w_2)+\rho(w_2)\rho(w_1)=2(w_1,w_2)$ applied to $\xi\ne0$ gives $(w_1,w_2)=0$. By Corollary \ref{cor:c7-Q}, for $(w_1,w_2)=0$ the component is a multiple of $\int\omega\,\big[(\iota_{w_1}A,(\rho(w_2)\xi,\lambda))+(\iota_{w_2}A,(\rho(w_1)\xi,\lambda))\big]$, which vanishes; the classical component is the same functional (proof of Corollary \ref{cor:c7-Q}).
\end{proof}

\begin{remark}[The pair entries of the holomorphic twist]
\label{rem:c11-twist}
The twisted one-loop differential is $\hbar Q^\xi_1(O)=\hbar L_1(O)^{(0)}+Q_1(O)-\tfrac12\hbar L_1(O)^{(2)}(\xi,\xi)-\cdots$ (Definition \ref{def:c8-twisted}), so its pair rule [BK, (3.7)--(3.8)] corresponds to the two-body kernels of $Q_1$ along $\xi$ and of the arity-two component along $(\xi,\xi)$. By Proposition \ref{prop:c7-select}, applied once for each argument, a component of arity $k$ along $\xi$ of an operator of Weyl weight $w$ and degree $d$ has representatives of Weyl weight $w+\tfrac k2$ and degree $d+1-k$. [BK, (3.51)] relates the holomorphic letters to the fields up to field redefinitions; with $\xi_\alpha{}^i=u_\alpha z^i$, $u'\in S_+$ with $u'^\alpha u_\alpha=1$, and the representatives $\gamma\leftrightarrow A_{ij}z^j$ (Weyl weight $1$, degree $0$), $\beta\leftrightarrow u'^\alpha\lambda_\alpha{}^i$ modulo $z^i$ $(\tfrac32,0)$, $\partial_{\dot\alpha}c\leftrightarrow\bar\lambda_{\dot\alpha i}z^i$ $(\tfrac32,0)$, $c\leftrightarrow c$ $(0,1)$, $b\leftrightarrow u'^\alpha F_{mn}(\sigma^m\bar\sigma^n)_\alpha{}^\beta u'_\beta$ $(2,0)$ or a component of $\lambda^*$ $(\tfrac52,-1)$, and $\partial_{\dot\alpha}\leftrightarrow u'^\alpha\sigma^m_{\alpha\dot\alpha}D_m$ $(1,0)$, the pair entries of [BK, (3.64)--(3.69)] lie in the following arities:

\begin{tabular}{l|c|c|c}
entry & Weyl weight and degree of the pair & of the output & arity\\
\hline
$b\,c$ [BK, (3.64)] & $(2,1)$ & $(3,0)$ & $2$\\
$b\,b$ [BK, (3.65)] & $(4,0)$ & $(\tfrac92,0)$ & $1$\\
$\beta\,\gamma$ [BK, (3.66)] & $(\tfrac52,0)$ & $(3,0)$ & $1$\\
$\beta\,\beta$ [BK, (3.67)] & $(3,0)$ & $(\tfrac72,0)$ & $1$\\
$b\,\gamma$ [BK, (3.68)] & $(3,0)$ & $(\tfrac72,0)$ & $1$\\
$\beta\,b$ [BK, (3.69)] & $(\tfrac72,0)$ & $(4,0)$ & $1$\\
\end{tabular}
\tabcaption{The arities of the one-loop pair entries of the holomorphic twist, with $b$ represented by the self-dual field strength; with $b$ represented by a component of $\lambda^*$ the entry $b\,c$ has arity $1$.\label{tab:c11-twist}}

The entry $\beta\gamma$, the Konishi entry, is the two-body kernel of $Q_1$ along $\xi$ on the pair $(\lambda|_{S_+\otimes W},A|_{W_{\mathbb R}^\perp})$ with two outputs $\lambda|_{S_-\otimes W^*}$: the coefficient of $\int\omega\,(\bar\lambda_{\dot\alpha i}z^i,\bar\lambda^{\dot\alpha}{}_jz^j)$ in $Q_1$ of $\int\omega\,(u'^\alpha\lambda_\alpha{}^i,A_{ij}z^j)$. The classical component of this pair along $\xi$ is a number $\kappa$ times a cubic functional of the scalars, and a finite part of the operators of Weyl weight $3$ that adds $y\int\omega\,(\bar\lambda_{\dot\alpha i}z^i,\bar\lambda^{\dot\alpha}{}_jz^j)$ to that cubic functional changes the entry by $-\kappa y$ (Proposition \ref{prop:c11-freedom} (iii)), the two functionals having the same Weyl weight and quantum numbers. The pair entries are therefore compared through the class $[Q^\xi_1]$ on the cohomology of $Q^\xi_0$ (Theorem \ref{thm:c8-twist}), transported to the holomorphic letters through [ESW, Thm 10.9] (remark after Theorem \ref{thm:c8-twist}). Proposition \ref{prop:c11-plane} is the exact untwisted form, for two letters, of [BK, (3.77)], $Q_1(\gamma^I\gamma^J)=0$: for $\rho(w_a)\xi=0$, $\iota_{w_a}A$ is a combination of the components $A_{ij}z^j$, the letters $\gamma$.
\end{remark}

\begin{remark}[The data that complete the one-loop components]
\label{rem:c11-next}
The two-body kernels of $Q_1$ along $\xi\in\Sigma$ are the cycles of Proposition \ref{prop:c10-arity1} with the odd line of Lemma \ref{lem:c11-lines} (i) or with the special edge and the vertex $S^{(1)}(\xi)$ on a leg; their evaluation by the simplex formulas of Theorem \ref{thm:c11-simplex}, in the scheme of Sections 11.3--11.5 and with the signs of Definition \ref{def:c11-sp}, is the next computation: first on the Konishi pair of Remark \ref{rem:c11-twist}, on the pairs of Table \ref{tab:c11-twist} and on the scalar bilinears of Corollary \ref{cor:c7-Q}. With them, $S_1$ follows by Proposition \ref{prop:c11-S} up to the term linear in $4\eta$, and $K_1$ by Corollary \ref{cor:c11-K} with the radial Gram matrix of the two-point graphs of Definition \ref{def:c8-2pt}. The cubic part of \eqref{eq:c9-arity1}, with the cubic part of the composite part from the composite loops of $I$ with three legs, is the consistency condition of Proposition \ref{prop:c11-theta-sector}, and the one-loop quantisation through arity two along $(\Delta,P_b)$ is the hypothesis of Theorem \ref{thm:c11-arity2}. The pole identity and the logarithmic covariance of Theorem \ref{thm:c11-local} (iii) are identities for the composite cycles through the pairs and the arity-one counterterms, tabulated with the derivatives of Table \ref{tab:c11-counts}; the gauge completion of the families couples the sector coefficients through the three-field parts.
\end{remark}

\begin{remark}[Two-component forms]
\label{rem:c11-twocomp}
With the components of Definition \ref{def:c8-twocomp}:
\begin{itemize}
\item Letters. $A|_V$, $A|_{W_{\mathbb R}^\perp}$, $\lambda|_{S_+\otimes W}$, $\lambda|_{S_-\otimes W^*}$ are $A_m$, $A_{ij}$ (the $\mathbf 6$ of $SU(4)_R$), $\lambda_\alpha{}^i$ (undotted, $\mathbf 4$), $\bar\lambda_{\dot\alpha i}$ (dotted, $\bar{\mathbf 4}$), with the antifields $A^*_m$, $A^*_{ij}$, $\lambda^{*\alpha}{}_i$, $\bar\lambda^{*\dot\alpha i}$, $c^*$ and the degrees $1,0,0,0,0,-1,-1,-1,-1,-2$ for $c,A_m,A_{ij},\lambda_\alpha{}^i,\bar\lambda_{\dot\alpha i},A^*_m,A^*_{ij},\lambda^{*\alpha}{}_i,\bar\lambda^{*\dot\alpha i},c^*$; $(\lambda,\lambda^*)=(\lambda_\alpha{}^i,\lambda^{*\alpha}{}_i)-(\bar\lambda_{\dot\alpha i},\bar\lambda^{*\dot\alpha i})$ and $(u,w)=\tfrac14u^{ij}w_{ij}$ on $W_{\mathbb R}^\perp$.
\item Composite kernels. The internal vertex along $\Delta$ is $\int\big((2E+2)A_m,A^*_m\big)=2\int x_n(\partial_nA_m,A^*_m)+2\int(A_m,A^*_m)$ with $E=x_n\partial_n$, along $P_b$ it is $\int(b_n\partial_nA_m,A^*_m)$; it sits between $\langle A_mA_n\rangle=2(-,-)^{-1}\delta_{mn}$ and $\langle A^*_m\,c\rangle=-(-,-)^{-1}ip_m$, each times $e^{-\varepsilon|p|^2}|p|^{-2}$ at $T=\infty$, and the insertion factors are $2$, $2i(l+q_c)_n$ and $ib_n(l+q_c)_n$. The special edge carries $e^{-\varepsilon|p|^2}-1$ times $(-,-)^{-1}$ and the inverse fibre pairing of $(A_m,A^*_m)$, $\tfrac14(A^{ij},A^*_{ij})$, $(\lambda,\lambda^*)$ and $(c,c^*)$: on $(\lambda_\alpha{}^i,\lambda^{*\beta}{}_j)$ the factor $-|l+q_\sigma|^2\delta_\alpha^\beta\delta^i_j(-,-)^{-1}$, on $(\bar\lambda_{\dot\alpha i},\bar\lambda^{*\dot\beta j})$ the opposite sign.
\item Bilinears. For $O=\tfrac1{16}\int\omega\,(w_1^{ij}A_{ij},w_2^{kl}A_{kl})$ the composite pair kernel along $\Delta$ is
\[
\Big[-\tfrac13\big((2E+2)A_{ij}\big)A_{kl}-\tfrac13A_{ij}\big((2E+2)A_{kl}\big)-\tfrac43A_{ij}A_{kl}\Big]\ \text{units},
\]
and along $P_b$ it is $-\tfrac13$ units times $b_m\partial_m(A_{ij}A_{kl})$, with the same contraction.
\item Two-site density. On adjacent pairs $(A_{ij},A_{kl})$, $(Z,\psi)$, $(\psi,Z)$, $(\psi,\psi)$ of a single trace the composite pair kernels give $-\tfrac43$ units (translation-invariant part along $\Delta$) and $-\tfrac16$ units times $b_m\partial_m$ (along $P_b$) on the same word, and $0$ on the exchanged word.
\item Free directions. The sector parts of $S^{(1)}(\Delta)$ are $\tfrac14\int\big((2E+2)A^{ij},A^*_{ij}\big)$, $-\int\big[\big((2E+3)\lambda_\alpha{}^i,\lambda^{*\alpha}{}_i\big)-\big((2E+3)\bar\lambda_{\dot\alpha i},\bar\lambda^{*\dot\alpha i}\big)\big]$ and $\int\big((2E+2)A_m,A^*_m\big)-\int\big(2E(c),c^*\big)$; those of $S^{(1)}(P_b)$ are $\tfrac14\int(b_m\partial_mA^{ij},A^*_{ij})$, $-\int\big[(b_m\partial_m\lambda_\alpha{}^i,\lambda^{*\alpha}{}_i)-(b_m\partial_m\bar\lambda_{\dot\alpha i},\bar\lambda^{*\dot\alpha i})\big]$ and $\int(b_m\partial_mA_n,A^*_n)-\int(b_m\partial_mc,c^*)$. The exact directions are $Q\int x_m(A_m,c^*)$ and $Q\int(\rho(x)\lambda^*,\lambda^*)=-2i\,Q\int x_m\big(\lambda^{*\alpha}{}_i,\sigma^m_{\alpha\dot\alpha}\bar\lambda^{*\dot\alpha i}\big)$. The further class of Proposition \ref{prop:c11-tc0} without the reflection is a multiple of $\int\big[(\lambda_\alpha{}^i,\lambda^{*\alpha}{}_i)+(\bar\lambda_{\dot\alpha i},\bar\lambda^{*\dot\alpha i})\big]$, the rotation $\lambda_\alpha{}^i\mapsto e^{i\varphi}\lambda_\alpha{}^i$, $\bar\lambda_{\dot\alpha i}\mapsto e^{-i\varphi}\bar\lambda_{\dot\alpha i}$ of the free gaugino.
\item Letter count. It counts $2$ for each letter $A_m$, $A_{ij}$, $\lambda_\alpha{}^i$, $\bar\lambda_{\dot\alpha i}$ and $1$ for each derivative $\partial_m$ on it; (b) reads: the composite pair kernels along $P_b$ sum to $-\tfrac13$ units times $b_m\partial_m$ on cohomology.
\item Families and spectra. $n_m=\sigma^m_{1\dot1}$, $D_n=n_m(\partial_m+[A_m,-])$, $\psi=\lambda_1{}^i$, and $F_{11}$ the self-dual field strength with both undotted indices $1$. In Srednicki's normalisation, $A_m=-igA^a_mT^a$, $A_I=-igA^a_IT^a$, $\lambda=-\tfrac{ig}{\sqrt2}\lambda^aT^a$, $\operatorname{tr}(T^aT^b)=\tfrac12\delta^{ab}$, the unit is $-\frac{g^2N}{32\pi^2}$ and $-\tfrac13$ units is $\frac{g^2N}{96\pi^2}$. On the families $D_1=-8$ units times the eigenvalues $E$ of $\sum_lH_{l,l+1}$, which is $\frac{g^2N}{4\pi^2}E=2\gamma$ with the one-loop anomalous dimension $\gamma=\frac{g^2N}{8\pi^2}E$: $\gamma=\frac{g^2N}{2\pi^2}h(S)$ on twist two, $\gamma=\frac{3g^2N}{4\pi^2}$ for $\operatorname{tr}(F_{11}F_{11})$ and the trace bilinear, the Konishi multiplet.
\item Odd generators. $Q_1(O)=\xi_\alpha{}^i\,Q_1(O)^\alpha{}_i+\bar\xi_{\dot\alpha i}\,\bar Q_1(O)^{\dot\alpha i}$ and $S_1(O)=\eta^\alpha{}_i\,S_1(O)_\alpha{}^i+\bar\eta^{\dot\alpha i}\,\bar S_1(O)_{\dot\alpha i}$ (Definition \ref{def:c8-components}); the Killing spinor of $\eta$ has the components $(\rho(x)\eta)_\alpha{}^i=-i\,x_m\sigma^m_{\alpha\dot\alpha}\bar\eta^{\dot\alpha i}$ and $(\rho(x)\eta)_{\dot\alpha i}=i\,x_m\eta^\alpha{}_i\sigma^m_{\alpha\dot\alpha}$ \eqref{eq:c8-pairing}, so on a family with $P_1\cong0$ the position part along $b$ of $S_1$ along $\eta$ is, on cohomology, $Q_1$ along $\xi_\alpha{}^i=-i\,b_m\sigma^m_{\alpha\dot\alpha}\bar\eta^{\dot\alpha i}$, $\bar\xi_{\dot\alpha i}=i\,b_m\eta^\alpha{}_i\sigma^m_{\alpha\dot\alpha}$. For $\xi_\alpha{}^i=u_\alpha z^i$, $(\rho(w)\xi)^\alpha{}_j=\varepsilon^{\alpha\beta}w_{ji}u_\beta z^i$, so $\rho(w)\xi=0$ exactly when $w_{ij}z^j=0$, and then $\iota_wA=\tfrac14w^{ij}A_{ij}$ is a combination of the components $A_{ij}z^j$: Proposition \ref{prop:c11-plane} is the vanishing of $Q_1$ on $\int\omega\,(m_1^iA_{ij}z^j,m_2^kA_{kl}z^l)$, $m_1,m_2\in W$.
\item Weinberg. On the Dirac column $(\lambda_\alpha{}^i,\bar\lambda^{\dot\alpha}{}_i)$, $\rho(v)=\gamma_W^\mu v_\mu$ gives Weinberg's form of each line, as in Remark \ref{rem:c10-twocomp}.
\end{itemize}
\end{remark}
